\section{Standar Operasional Prosedur dan Kerangka Kelembagaan}
\label{sec:pelaksanaan}

\subsection{Standar Operasional Prosedur Pelaksanaan}

Urutan waktu menentukan keberhasilan program. Pembangkit pengganti membutuhkan beberapa tahun untuk dibangun, sehingga keputusan pensiun dini harus diambil paling lambat tiga tahun sebelum tanggal penutupan. Tabel~\ref{tab:sop} merinci prosedur untuk setiap PLTU kandidat.

\begin{table}[H]
\centering
\caption{SOP pasangan kuota PLTS atap dan pensiun dini per PLTU}
\label{tab:sop}
\small
\begin{tabularx}{\textwidth}{L{2.6cm} Y L{2.6cm} L{2.4cm}}
\toprule
Tahap & Prosedur & Penanggung jawab & Waktu \\
\midrule
1. Penyaringan kandidat & Menjalankan tab Kandidat PLTU di PowerShift dengan data kontrak PLN, mengisi lembar penilaian (Lampiran~\ref{lamp:lembar}), menetapkan daftar urut & Satgas, PLN & Fase 1 (2026-2027), diperbarui tiap tahun \\
\addlinespace
2. Penetapan paket & Menyusun skenario di PowerShift berisi PLTU, tahun pensiun, kuota PLTS atap, dan pendanaan, lalu memeriksa empat syarat penutupan & Satgas, KESDM & $Y-3$ \\
\addlinespace
3. Kesepakatan kompensasi & Negosiasi dengan IPP berbasis keuntungan yang hilang, pengikatan dana luar paling sedikit 80\% & Kemenkeu, PLN, mitra ETM dan JETP & $Y-3$ sampai $Y-2$ \\
\addlinespace
4. Penyiapan pengganti & Menjadwalkan pembangkit pengganti dalam RUPTL dan memulai pembangunan & PLN & $Y-3$ sampai $Y$ \\
\addlinespace
5. Kuota dan kewajiban PLTS atap & Menetapkan kuota sistem, menjalankan kewajiban gedung pemerintah dan BUMN, menyalurkan insentif industri & KESDM, Kemenkeu, Kementerian BUMN & Tahunan sejak 2028 \\
\addlinespace
6. Transisi pekerja & Mendata pekerja, memulai pelatihan ulang dan perlindungan sosial & Kemenaker, Pemda & $Y-1$ \\
\addlinespace
7. Penutupan dan pemantauan & Menutup PLTU setelah syarat terpenuhi, memperbarui skenario di PowerShift dengan realisasi, melaporkan indikator & PLN, Satgas & $Y$ dan tahunan \\
\bottomrule
\end{tabularx}
\sumber{disusun kelompok. $Y$ adalah tahun penutupan PLTU}
\end{table}

\subsection{Kerangka Kelembagaan}

Pembagian peran mengikuti matriks aktor laporan kelompok \cite{L-LAPORAN-K17} dan tata kelola transisi yang dibahas dalam literatur \cite{P-RESOSUDARMO-2023}.

\textbf{Tingkat koordinasi.} Satgas Transisi Energi lintas kementerian menetapkan daftar urut kandidat, memeriksa syarat penutupan, dan meninjau indikator di akhir setiap fase.

\textbf{Tingkat kementerian.} Kementerian ESDM menetapkan kuota berpasangan, kewajiban PLTS atap, dan petunjuk teknis pensiun dini. Kementerian Keuangan mengelola insentif industri, anggaran PLTS atap gedung pemerintah, dan pengikatan dana luar. Kementerian BUMN memastikan kewajiban PLTS atap di gedung BUMN. KLHK mengawasi pemulihan lahan bekas PLTU.

\textbf{Tingkat operator.} PT PLN (Persero) menyusun kajian pensiun dini per PLTU, memasukkan data kontrak ke PowerShift, menyiapkan pembangkit pengganti, dan memperkuat jaringan distribusi untuk PLTS atap.

\textbf{Tingkat daerah dan pendukung.} Kemenaker dan pemerintah daerah menjalankan pelatihan ulang. Sektor swasta dan perbankan membiayai PLTS atap industri. Perguruan tinggi memelihara PowerShift dan memverifikasi parameter sebagai barang publik.

\subsection{Alur Data dan Pelaporan}

Data mengalir berjenjang melalui PowerShift. PLN memasukkan parameter kontrak dan realisasi pembangkitan setiap tahun. KESDM memasukkan realisasi PLTS atap per provinsi. Satgas menjalankan skenario realisasi dan membandingkannya dengan skenario roadmap. Hasil agregat, yaitu neraca per pihak, indikator, dan status syarat penutupan, dibuka kepada publik sebagai tautan skenario, sedangkan data kontrak yang rahasia tetap berada di PLN.

\subsection{Peta Jalan 2026-2040}

\begin{table}[H]
\centering
\caption{Peta jalan roadmap dan keluaran terukur}
\label{tab:fase}
\small
\begin{tabularx}{\textwidth}{L{1.8cm} Y L{4.4cm}}
\toprule
Fase & Langkah utama & Keluaran simulasi akhir fase \\
\midrule
Fase 1 \newline 2026-2027 & Membentuk Satgas, menerbitkan peraturan kuota berpasangan (Lampiran~\ref{lamp:permen}), menetapkan daftar urut kandidat, memulai negosiasi Cirebon-1 & PLTS atap 910 MWp. Belum ada perubahan neraca \\
\addlinespace
Fase 2 \newline 2028-2030 & Kewajiban PLTS atap 200 MW/tahun dan insentif industri 15\% berlaku. Cirebon-1 pensiun 2030 setelah 660 MW pengganti siap & PLTS atap 2.307 MWp. PLTU pensiun 660 MW. Neraca negara -Rp7,4 T. 827 kematian dini dihindari \\
\addlinespace
Fase 3 \newline 2031-2035 & Tanjung Jati B 1-2 pensiun 2031, Cilacap Phase I 2033, Celukan Bawang 2035. Penguatan jaringan dan penyimpanan energi & PLTS atap 4.764 MWp. PLTU pensiun 3.006 MW. Neraca negara -Rp0,6 T. Dana luar Rp10,1 T \\
\addlinespace
Fase 4 \newline 2036-2040 & Evaluasi menyeluruh, penyaringan ulang kandidat termasuk PLTU milik PLN, pemulihan ekonomi daerah & PLTS atap 7.220 MWp. Kumulatif 9.572 kematian dini dan 35,4 Mt CO$_2$ dihindari \\
\bottomrule
\end{tabularx}
\sumber{roadmap laporan kelompok \cite{L-LAPORAN-K17}, keluaran dari simulasi PowerShift}
\end{table}
